% LaTeX sample document
% This is a comment

\documentclass[12pt,a4paper]{article}

% Package imports
\usepackage[utf8]{inputenc}
\usepackage{amsmath}
\usepackage{amssymb}
\usepackage{graphicx}
\usepackage{hyperref}
\usepackage{listings}
\usepackage{xcolor}

% Document metadata
\title{Sample LaTeX Document}
\author{John Doe \and Jane Smith}
\date{\today}

% Custom commands
\newcommand{\R}{\mathbb{R}}
\newcommand{\N}{\mathbb{N}}
\newcommand{\Z}{\mathbb{Z}}
\newcommand{\deriv}[2]{\frac{\mathrm{d}#1}{\mathrm{d}#2}}
\newcommand{\pderiv}[2]{\frac{\partial #1}{\partial #2}}

\begin{document}

\maketitle

\begin{abstract}
This is a sample LaTeX document demonstrating various features including text formatting, mathematical equations, figures, tables, and references. The document serves as a comprehensive example for learning LaTeX syntax and structure.
\end{abstract}

\tableofcontents

\section{Introduction}

This is a sample \LaTeX{} document. LaTeX is a typesetting system commonly used for scientific and technical documents. It provides excellent support for mathematical notation and automatic numbering of sections, equations, figures, and tables.

\subsection{Text Formatting}

LaTeX supports various text formatting options:
\begin{itemize}
    \item \textbf{Bold text} using \texttt{\textbackslash textbf}
    \item \textit{Italic text} using \texttt{\textbackslash textit}
    \item \underline{Underlined text} using \texttt{\textbackslash underline}
    \item \emph{Emphasized text} using \texttt{\textbackslash emph}
    \item \texttt{Monospace text} using \texttt{\textbackslash texttt}
\end{itemize}

\subsection{Lists}

\begin{enumerate}
    \item First item
    \item Second item
    \item Third item with sub-items:
    \begin{enumerate}
        \item Sub-item a
        \item Sub-item b
    \end{enumerate}
\end{enumerate}

\section{Mathematical Expressions}

\subsection{Inline Mathematics}

Inline math can be written using dollar signs: $E = mc^2$, or the equation $a^2 + b^2 = c^2$ for right triangles.

\subsection{Display Mathematics}

The quadratic formula is:
\[
x = \frac{-b \pm \sqrt{b^2 - 4ac}}{2a}
\]

An integral example:
\[
\int_{0}^{\infty} e^{-x^2} \, dx = \frac{\sqrt{\pi}}{2}
\]

A summation example:
\[
\sum_{n=1}^{\infty} \frac{1}{n^2} = \frac{\pi^2}{6}
\]

\subsection{Equations with Numbers}

\begin{equation}
\label{eq:euler}
e^{i\pi} + 1 = 0
\end{equation}

Equation \ref{eq:euler} is Euler's identity, which connects five fundamental mathematical constants.

\begin{equation}
\label{eq:fourier}
f(x) = \frac{a_0}{2} + \sum_{n=1}^{\infty} \left( a_n \cos(nx) + b_n \sin(nx) \right)
\end{equation}

\subsection{Matrices}

A matrix can be written as:
\[
A = \begin{pmatrix}
a_{11} & a_{12} & a_{13} \\
a_{21} & a_{22} & a_{23} \\
a_{31} & a_{32} & a_{33}
\end{pmatrix}
\]

The determinant:
\[
\det(A) = \begin{vmatrix}
a & b \\
c & d
\end{vmatrix} = ad - bc
\]

\section{Figures and Tables}

\subsection{Figures}

\begin{figure}[h]
\centering
% \includegraphics[width=0.8\textwidth]{example-image}
\caption{An example figure}
\label{fig:example}
\end{figure}

Figure \ref{fig:example} shows an example image.

\subsection{Tables}

\begin{table}[h]
\centering
\begin{tabular}{|c|c|c|}
\hline
\textbf{Column 1} & \textbf{Column 2} & \textbf{Column 3} \\
\hline
Row 1, Col 1 & Row 1, Col 2 & Row 1, Col 3 \\
Row 2, Col 1 & Row 2, Col 2 & Row 2, Col 3 \\
Row 3, Col 1 & Row 3, Col 2 & Row 3, Col 3 \\
\hline
\end{tabular}
\caption{An example table}
\label{tab:example}
\end{table}

Table \ref{tab:example} demonstrates a basic table structure.

\section{Code Listings}

\begin{lstlisting}[language=Python]
def fibonacci(n):
    if n <= 1:
        return n
    return fibonacci(n-1) + fibonacci(n-2)

# Calculate the 10th Fibonacci number
result = fibonacci(10)
print(f"The 10th Fibonacci number is {result}")
\end{lstlisting}

\section{References and Citations}

This document cites the work of \cite{knuth1984} and \cite{lamport1994}.

\subsection{Cross-references}

\begin{itemize}
    \item See Equation \ref{eq:euler} on page \pageref{eq:euler}
    \item See Figure \ref{fig:example} on page \pageref{fig:example}
    \item See Table \ref{tab:example} on page \pageref{tab:example}
\end{itemize}

\section{Theorems and Proofs}

\begin{theorem}[Pythagorean Theorem]
For a right triangle with legs of length $a$ and $b$ and hypotenuse $c$:
\[
a^2 + b^2 = c^2
\]
\end{theorem}

\begin{proof}
Consider a right triangle with legs $a$ and $b$ and hypotenuse $c$. By constructing a square with side $a+b$ and arranging four copies of the triangle, we can prove that the area of the inner square equals $c^2$, which equals $a^2 + b^2$.
\end{proof}

\section{Conclusion}

This document has demonstrated various LaTeX features including:
\begin{itemize}
    \item Document structure and metadata
    \item Text formatting and lists
    \item Mathematical expressions
    \item Figures and tables
    \item Code listings
    \item References and citations
\end{itemize}

\bibliographystyle{plain}
\bibliography{references}

\appendix

\section{Appendix}

Additional content can be placed in appendices.

\end{document}
